\section{Aufbau}
Es werden zwei Informationsstrategien entwickelt, ein Dialogbasiertes System und 
eine Website die als Multi-Step-Funnel gestaltet wird.

\subsection{Dialogsystem}

Das Dialogbasierte System
wird als KI-Chatbot umgesetzt. Der Chatbot wird so gestaltet, dass er zum einen,
vom Nutzer gestellte Fragen über die beworbene App beantwortet, darüber hinaus
aber auch versucht Interesse bei seinem Gesprächspartner zu wecken, und diesem
die Vorteile von Jobbing anhand der gestellten Fragen zeigt.
Um eine Evaluierung des bestmöglichen Chatbots zu ermöglichen wird der Chatbot
so entwickelt, dass der Chatbot auf verschiedene Arten mit dem Nutzer kommuniziert.

Um eine neutrale Auswertung zu ermöglichen wird dem KI-Assistent zu beginn eine
zufällige Kommunikationsstrategie zugewiesen. Der Chatbot behält diese Strategie
über die ganze Konversation hinweg. Die Unterschiedlichen Varianten werden dann
Anhand der erhobenen Daten verglichen und es wird der effektivste ausgewählt. 
Die beste der Varianten wird dann für den Vergleich mit dem Multi-Step-Funnel
herangezogen, da nur so ein aussagekräftiger Vergleich zwischen den Konzepten
möglich ist. Wenn man beispielsweise eine schlechte Chatbot Variante mit einem
guten Funnel vergleicht, kann zwar eine Aussage über genau diese Varianten 
getroffen werden, die Daten können jedoch nicht für eine generelle Aussage
über die grundlegenden Konzepte verwendet werden, da schlechte Umsetzungen 
guter Konzepte im Marketing nicht verwendet werden.

\textit{hier noch Evaluierung warum genau diese Varianten getestet werden}
Es werden der Einfluss von Tonalität und Personalisierung auf das Nutzerverhalten
untersucht. Um die Faktoren unabhängig testen zu können und so die beste Kombination
ersichtlich zu machen werden vier verschiedene Persönlichkeitsvarianten definiert.
Folgende Kommunikationsstrategien sind möglich:

\begin{itemize}
    \item formeller Ton und wenig Personalisierung
    \item informeller Ton und wenig Personalisierung
    \item formeller Ton und viel Personalisierung
    \item informeller Ton und viel Personalisierung
\end{itemize}

\subsection{Funnel}

Der Multi-Step-Funnel beantwortet auch allgemeine Fragen über die App in Form
von informativem Text und versucht durch interaktive Benutzeroberfläche den 
Nutzer interessiert zu halten und Informationsüberflutung zu vermeiden.
Während des Durchlaufs werden dem Nutzer die Vorteile erklärt und zum Schluss
erhält der Verbraucher die Möglichkeit zum Download.


\section{Zielgruppen}

Eine der wichtigsten Herausforderungen, in der Gestaltung der digitalen Nutzerführung, besteht
darin, die angebotenen Inhalte so zu gestalten, dass sie für möglichst viele
unterschiedliche Nutzergruppen relevant sind und somit ansprechend sind.
Personalisierung und gezielte Ansprache spezieller Nutzergruppen gelten als wirksam,
um Wahrnehmung der Informationsangebote zu steigern und die Vermittlung zu optimieren. In
der Marketing- und Werbeforschung ist es immer mehr belegt worden, dass personalisierte
Botschaften nicht nur intensiver wahrgenommen werden, sondern auch deutlich bessere
Effekte in der Informationswahrnehmung haben als nicht personalisierte verallgemeinerte Botschaften.

Eine aktuelle Meta-Analyse experimenteller Studien zu personalisierter Werbung kommt
zu dem Ergebnis, dass personalisierte Werbung im Vergleich zu generischen Varianten in der
Regel wirksamer ist, um Einstellungen zu beeinflussen und die Intentionen der Wahrnehmenden
zu verändern.
Speziell personalisierte Werbebotschaften führen zu günstigeren
Konsumenten-Einstellungen und stärkeren Verhaltensabsichten als nicht personalisierte.
Die Autoren erklären diesen Effekt in erster Linie damit, dass Personalisierung die
subjektiv empfundene Relevanz der Botschaft erhöht und damit auch die persuasive
Wirksamkeit steigert.
Die Ergebnisse Der Meta-Analyse zeigen sehr deutlich, dass Personalisierung im digitalen
interaktiven Marketing sinnvoll ist und vorteilhaft eingesetzt werden kann.
\cite{Yeo2025metaPersonalizedAdvertising}

Aufgrund dessen ist es sinnvoll, sich bei der Gestaltung des Funnels zu überlegen auf Welche
Zielgruppen die Gestaltung angepasst werden soll, um eine möglichst hohe Effektivität bei
der gewünschten Zielgruppe zu erreichen. Im Grunde kann in dieser Arbeit zwischen zwei 
Grundlegenden Zielgruppen unterschieden werden: Gastronom:innen und Jobsuchende.

Ob diese Personalisierung auch bei der verbalen Kommunikation mit einem Chatbot vorteilhaft
ist oder ob sie hier möglicherweise sogar als störend empfunden wird, ist eine wichtige
Frage, die durch die Analyse des Nutzerverhaltens in verschiedenen Kommunikationsstrategien
untersucht und beantwortet werden soll.

\subsection{Gastronom:innen}

\subsection{Jobsuchende}

\section{Funktionale Anforderungen}

\section{Sonstige Anforderungen}

\section{Use-Cases}

\section{Funnel Konzept}

\subsection{Klassisch}

\subsection{Chatbot basiert}

\section{Messungen}
